\documentclass[10pt, letterpaper]{article}

\usepackage[margin=1in]{geometry}
\usepackage[T1]{fontenc}
\usepackage{lmodern}
\usepackage{amsmath, amssymb, amsthm}
\usepackage{booktabs}
\usepackage{graphicx}
\usepackage{xcolor}
\usepackage{hyperref}
\usepackage{enumitem}
\usepackage{caption, subcaption}
\usepackage{listings}
% \usepackage{stmaryrd}  % not available in basic TeX Live

% Theorem environments
\newtheorem{theorem}{Theorem}[section]
\newtheorem{lemma}[theorem]{Lemma}
\newtheorem{proposition}[theorem]{Proposition}
\newtheorem{corollary}[theorem]{Corollary}
\theoremstyle{definition}
\newtheorem{definition}[theorem]{Definition}
\newtheorem{example}[theorem]{Example}
\theoremstyle{remark}
\newtheorem{remark}[theorem]{Remark}

% Lean code formatting
\lstdefinelanguage{Lean}{
  morekeywords={def, theorem, lemma, axiom, variable, inductive, structure,
    class, instance, where, match, with, by, exact, intro, apply, simp,
    rfl, fun, let, have, show, calc, if, then, else, do, return, sorry,
    namespace, end, open, import, section, Type, Prop, True, False,
    forall, exists, private, set, obtain, push_neg, unfold, constructor,
    rintro, refine, subst, rw, cases, induction, congr, ext},
  sensitive=true,
  morecomment=[l]{--},
  morecomment=[n]{/-}{-/},
  morestring=[b]",
  literate={→}{$\to$}1 {←}{$\leftarrow$}1 {↔}{$\leftrightarrow$}1
           {∀}{$\forall$}1 {∃}{$\exists$}1 {¬}{$\lnot$}1
           {∧}{$\land$}1 {∨}{$\lor$}1 {⟨}{$\langle$}1 {⟩}{$\rangle$}1
           {₁}{$_1$}1 {₂}{$_2$}1 {₀}{$_0$}1
           {⊊}{$\subsetneq$}1,
}
\lstset{
  language=Lean,
  basicstyle=\ttfamily\small,
  keywordstyle=\bfseries\color{blue!70!black},
  commentstyle=\itshape\color{green!50!black},
  stringstyle=\color{red!60!black},
  breaklines=true,
  columns=flexible,
  frame=single,
  framesep=4pt,
  xleftmargin=8pt,
  xrightmargin=8pt,
  aboveskip=8pt,
  belowskip=8pt,
}

% Notation
\newcommand{\Sach}{\mathsf{S}}
\newcommand{\World}{\mathsf{World}}
\newcommand{\Prop}{\mathsf{Prop}}
\newcommand{\eval}{\mathsf{eval}}
\newcommand{\expresses}{\mathsf{expresses}}
\newcommand{\structEq}{\sim_{\mathrm{struct}}}
\newcommand{\formEq}{\sim_{\mathrm{form}}}
\newcommand{\semEq}{\sim_{\mathrm{sem}}}
\newcommand{\Lean}{\textsc{Lean}}

\title{What \Lean{} Cannot Say:\\
A Machine-Checked Analysis of Wittgenstein's \emph{Tractatus}}
\author{Robb Walters}
\date{April 2026}

\begin{document}
\maketitle

\begin{abstract}
We present a machine-checked reconstruction of the semantic core
of \emph{Tractatus Logico-Philosophicus} in the \Lean{} proof
assistant.  The formalization models objects, states of affairs,
worlds, and propositions, and proves truth-functional
compositionality, functional completeness, and bivalence.

Beyond encoding, the system is parameterized over world models,
allowing a precise separation between (i)~structural invariants
of the propositional syntax, (ii)~model-dependent assumptions
such as the independence of atomic facts, and (iii)~formal limits
of expressibility.  We show that truth-functional compositionality
holds in all models, while independence characterizes a specific
``free'' model and fails in constrained models.

To analyze the relation between object language and metalanguage,
we introduce an \texttt{expresses} relation linking propositions
to meta-level predicates.  Using this, we prove that any
proposition expressing a world-independent property must be a
tautology or contradiction.  Consequently, nontrivial propositions
cannot express world-independent truths.

We further distinguish three notions of
equivalence---syntactic identity, logical-form equivalence up to
permutation, and semantic equivalence---and prove their strict
separation.

These results yield a formal decomposition of the \emph{Tractatus}
into invariants, assumptions, and limits, and provide a precise
characterization of the saying/showing distinction as a theorem
about expressibility in truth-functional languages.
\end{abstract}

% ================================================================
\section{Introduction}
\label{sec:intro}
% ================================================================

Wittgenstein's \emph{Tractatus Logico-Philosophicus}~\cite{wittgenstein1921}
proposes that the world consists of atomic facts (Sachverhalte),
that propositions are truth-functions of elementary propositions,
and that the limits of language are the limits of the world.
The text is terse, aphoristic, and deliberately resistant to
formalization---yet its semantic core is a precise theory of
truth-functional composition that invites formal reconstruction.

Several authors have pursued this invitation.
Lokhorst~\cite{lokhorst1988} provides a set-theoretic
reconstruction covering ontology, semantics, and propositional
attitudes;
Miller~\cite{miller1995} develops a Tractarian semantics for
predicate logic;
Weiss~\cite{weiss2017} reconstructs the logical system with
$\Pi^1_1$-completeness results;
Stokhof~\cite{stokhof2008} traces the \emph{Tractatus}'s
influence on formal semantics.
These reconstructions clarify the \emph{Tractatus} as a semantic
system but are carried out on paper, without machine verification.
Separately, the proof-assistant community has formalized
philosophical arguments~\cite{benzmuller2017,fitelson2007} and
even complete metaphysical systems~\cite{scherf2025}, but a
Tractarian formalization has not appeared.

This paper contributes a machine-checked reconstruction in
\Lean~4~\cite{moura2021} that goes beyond encoding to yield new
analytical results.  Our central claim is:

\begin{quote}
\emph{The saying/showing distinction of the \emph{Tractatus} can
be reformulated as a theorem: nontrivial propositions cannot
express world-independent truths.}
\end{quote}

\noindent
This is not a paraphrase of Wittgenstein but a provable
consequence of the typed structure of the formalization.
The key device is an \texttt{expresses} relation (Definition~\ref{def:expresses})
that bridges the object language (propositions of type
$\mathsf{Proposition}\;\Sach$) and the metalanguage (terms of
type $\Prop$), making the object/meta boundary visible at the
type level.

More broadly, by parameterizing the formalization over a general
\texttt{WorldModel} structure, we show that the claims of the
\emph{Tractatus} decompose into three classes:
\begin{enumerate}[nosep]
  \item \textbf{Structural invariants} that hold in every world
    model---compositionality, bivalence, De Morgan duality.
  \item \textbf{Model-dependent assumptions} that hold in the
    unconstrained ``free'' model but fail in constrained
    models---principally, the independence of atomic facts.
  \item \textbf{Formal limits} on expressibility---the collapse
    of world-independent content to tautology or contradiction,
    and the strict separation of syntactic, logical-form, and
    semantic equivalence.
\end{enumerate}

\noindent
The formalization comprises 1,517 lines of \Lean~4, proves 46
theorems with 0 \texttt{sorry} obligations and 1 deliberate
axiom (\texttt{axiom silence : True}, encoding TLP~7), and is
publicly available.\footnote{\url{https://leangenius.org/proof/tractatus-ontology}}

The paper is organized as follows.
Section~\ref{sec:formalization} presents the formalization:
types, evaluation, and the world-model abstraction.
Section~\ref{sec:invariants} establishes the structural
invariants.
Section~\ref{sec:assumptions} analyzes the model-dependent
assumptions and exhibits countermodels.
Section~\ref{sec:limits} proves the expressibility results and
the equivalence hierarchy.
Section~\ref{sec:related} discusses related work.
Section~\ref{sec:discussion} reflects on the philosophical
implications and limitations.

% ================================================================
\section{The Formalization}
\label{sec:formalization}
% ================================================================

We formalize the ontological skeleton of the \emph{Tractatus}
in \Lean~4.  The design choices are deliberate modeling
decisions, not exegetical claims about Wittgenstein's intent.
We state them explicitly so that the reader can assess the
reconstruction's scope and limitations.

% ----------------------------------------------------------------
\subsection{Objects and States of Affairs}
\label{sec:objects}
% ----------------------------------------------------------------

TLP~2.02: ``The object is simple.''  Objects are the substance
of the world (TLP~2.021).  We model them as an abstract type
\texttt{TractObject} with no internal structure, capturing
Wittgenstein's insistence on simplicity.

A state of affairs (Sachverhalt) is a possible combination of
objects (TLP~2.01).  Rather than encoding the combinatorial
structure of objects into states of affairs, we index states of
affairs by an abstract type~$\Sach$:

\begin{lstlisting}
variable (TractObject : Type)
variable (Sachverhalt : Type)
\end{lstlisting}

\noindent
This choice leaves the internal structure of states of affairs
unspecified---an interpretive decision that captures independence
while remaining agnostic about combinatorial form (cf.\
TLP~2.0141).

% ----------------------------------------------------------------
\subsection{Worlds}
\label{sec:worlds}
% ----------------------------------------------------------------

TLP~1: ``The world is everything that is the case.''
A world is determined by which states of affairs obtain.  We
model this as a predicate:

\begin{lstlisting}
def World := Sachverhalt → Prop
\end{lstlisting}

\noindent
Every function $\Sach \to \Prop$ counts as a world.  This gives
the full Boolean cube of truth-value assignments, making the
independence of atomic facts (TLP~2.061) hold by construction.
We return to this modeling choice and its consequences in
Section~\ref{sec:assumptions}.

% ----------------------------------------------------------------
\subsection{Propositions}
\label{sec:propositions}
% ----------------------------------------------------------------

TLP~5: ``The proposition is a truth-function of elementary
propositions.''  We define propositions as an inductive type
built from elementary propositions, negation, and conjunction:

\begin{lstlisting}
inductive Proposition (S : Type) where
  | elementary : S → Proposition S
  | neg        : Proposition S → Proposition S
  | conj       : Proposition S → Proposition S → Proposition S
\end{lstlisting}

\noindent
Disjunction, implication, biconditional, and NAND are defined as
derived operations (TLP~5.101, 5.5), confirming that $\{\lnot, \land\}$
is an adequate basis.

% ----------------------------------------------------------------
\subsection{Semantic Evaluation}
\label{sec:eval}
% ----------------------------------------------------------------

The core semantic function interprets propositions against a
world:

\begin{lstlisting}
def Proposition.eval (p : Proposition S) (w : World S) : Prop :=
  match p with
  | .elementary s => w s
  | .neg q        => ¬ (q.eval w)
  | .conj q r     => q.eval w ∧ r.eval w
\end{lstlisting}

\noindent
An elementary proposition $\mathtt{elementary}\;s$ is true in
world~$w$ exactly when $w\;s$ holds.  Negation and conjunction
lift to their classical counterparts.  We also provide a
\texttt{Bool}-valued evaluator \texttt{evalBool} for decidable
computation and prove its correctness against
\texttt{eval}.\footnote{Theorem \texttt{evalBool\_correct}:
\texttt{(p.evalBool w = true) $\leftrightarrow$ p.eval
(fun s $\Rightarrow$ w s = true)}.}

% ----------------------------------------------------------------
\subsection{Tautology, Contradiction, Nontriviality}
\label{sec:tautology}
% ----------------------------------------------------------------

A proposition is a \emph{tautology} if it holds in every world,
a \emph{contradiction} if it holds in none, and
\emph{nontrivial} if it holds in some world and fails in
another:

\begin{lstlisting}
def IsTautology (p : Proposition S) : Prop :=
  ∀ w : World S, p.eval w
def IsContradiction (p : Proposition S) : Prop :=
  ∀ w : World S, ¬ (p.eval w)
def Nontrivial (p : Proposition S) : Prop :=
  ∃ w₁ w₂ : World S, p.eval w₁ ∧ ¬ p.eval w₂
\end{lstlisting}

\noindent
We prove that nontriviality is equivalent to being neither a
tautology nor a contradiction
(Theorem~\texttt{nontrivial\_iff\_not\_taut\_not\_contra}).

% ----------------------------------------------------------------
\subsection{General World Models}
\label{sec:worldmodel}
% ----------------------------------------------------------------

The definitions above hardwire the world model: every function
$\Sach \to \Prop$ is a world.  In many applications---physics,
epistemic logic, constrained databases---not every assignment is
legitimate.  We introduce a \texttt{WorldModel} structure to
parameterize the semantics:

\begin{lstlisting}
structure WorldModel (S : Type) where
  W        : Type
  holds    : W → S → Prop
  nonempty : Nonempty W
\end{lstlisting}

\noindent
The \emph{free model} \texttt{freeModel} recovers the original
semantics: $W = \Sach \to \Prop$ and \texttt{holds} is function
application.  We define a general evaluator \texttt{evalM} and
prove that \texttt{evalM (freeModel S)} agrees with
\texttt{Proposition.eval} (Theorem~\texttt{evalM\_free\_eq\_eval}).

This abstraction is the key structural device of the
formalization.  It allows us to state precisely which results
depend on the choice of world model and which do not.

% ================================================================
\section{Structural Invariants}
\label{sec:invariants}
% ================================================================

The following results hold in \emph{every} world model---they
are consequences of the inductive structure of propositions and
classical logic alone, with no assumptions about which worlds
exist.

% ----------------------------------------------------------------
\subsection{Truth-Functional Compositionality}
\label{sec:compositionality}
% ----------------------------------------------------------------

\begin{theorem}[Compositionality, model-independent]
\label{thm:compositionality}
For any world model $M$, any proposition $p$, and any two
worlds $w_1, w_2 \in M.W$: if $M.\mathit{holds}\;w_1\;s
\leftrightarrow M.\mathit{holds}\;w_2\;s$ for every state of
affairs~$s$, then $\eval_M\;p\;w_1 \leftrightarrow
\eval_M\;p\;w_2$.
\end{theorem}

\noindent
The proof is by structural induction on propositions.
The elementary case follows from the hypothesis; negation and
conjunction lift via \texttt{Iff.not} and \texttt{Iff.and}.

\begin{lstlisting}
theorem truth_functional_compositionality_gen (p : Proposition S)
    (w₁ w₂ : M.W)
    (h : ∀ s : S, M.holds w₁ s ↔ M.holds w₂ s) :
    evalM M p w₁ ↔ evalM M p w₂ := by
  induction p with
  | elementary s => exact h s
  | neg q ih     => simp only [evalM]; exact ih.not
  | conj q r ihq ihr => simp only [evalM]; exact ihq.and ihr
\end{lstlisting}

\noindent
This establishes TLP~5 as a structural invariant: the
truth-value of a compound proposition is determined by its
elementary constituents regardless of which worlds the model
admits.

% ----------------------------------------------------------------
\subsection{Bivalence and Classical Tautologies}
\label{sec:bivalence}
% ----------------------------------------------------------------

Elementary bivalence (TLP~4.023) is an instance of classical
excluded middle:

\begin{theorem}[Bivalence]
For any state of affairs~$s$ and world~$w$: $w\;s \lor \lnot\;w\;s$.
\end{theorem}

\noindent
The standard tautologies follow: excluded middle ($p \lor
\lnot p$), double negation ($\lnot\lnot p \leftrightarrow p$),
De Morgan's laws, and modus ponens.  We verify each as a formal
theorem.

% ----------------------------------------------------------------
\subsection{NAND Functional Completeness}
\label{sec:nand}
% ----------------------------------------------------------------

TLP~5.5 anticipates the Sheffer stroke.  We prove that negation
and conjunction are both expressible via NAND:

\begin{theorem}[NAND completeness]
\label{thm:nand}
For any propositions $p$, $q$ and world $w$:
\begin{enumerate}[nosep]
  \item $\mathtt{nand}(p, p).\eval\;w \leftrightarrow
    (\lnot p).\eval\;w$
  \item $(\lnot\;\mathtt{nand}(p, q)).\eval\;w \leftrightarrow
    (p \land q).\eval\;w$
\end{enumerate}
\end{theorem}

\noindent
Since $\{\lnot, \land\}$ generates all truth-functions and NAND
generates $\{\lnot, \land\}$, NAND alone is functionally
complete---confirming TLP~5.5.

% ================================================================
\section{Model-Dependent Assumptions}
\label{sec:assumptions}
% ================================================================

Not every claim in the \emph{Tractatus} is a structural
invariant.  The independence of atomic facts---perhaps the most
famous ontological thesis of the text---is a property of a
specific world model.

% ----------------------------------------------------------------
\subsection{Independence in the Free Model}
\label{sec:independence}
% ----------------------------------------------------------------

We formalize independence as a typeclass:

\begin{lstlisting}
class IndependentWorlds (S : Type) where
  realizable : ∀ assignment : S → Prop,
    ∃ w : World S, ∀ s, w s ↔ assignment s
\end{lstlisting}

\noindent
The free model trivially satisfies this: every assignment
\emph{is} a world.

\begin{lstlisting}
instance : IndependentWorlds S :=
  ⟨fun a => ⟨a, fun _ => Iff.rfl⟩⟩
\end{lstlisting}

\noindent
Independence is thus not a theorem about logic but a consequence
of the modeling choice $\World\;\Sach = \Sach \to \Prop$.
This distinction---between what is proved and what is assumed by
the encoding---is precisely the kind of clarification that a
typed formalization forces.

% ----------------------------------------------------------------
\subsection{Failure in Constrained Models}
\label{sec:constrained}
% ----------------------------------------------------------------

To show that independence is genuinely contingent on the model,
we exhibit two countermodels.

\paragraph{Abstract countermodel.}
Consider two states of affairs $a, b$ with $a \neq b$, where
$b$ must co-occur with $a$.  Define
$\mathtt{ConstrainedWorld}\;\Sach\;a\;b = \{w : \Sach \to \Prop
\mid w\;a \to w\;b\}$.  Then:

\begin{theorem}[Independence fails in constrained models]
\label{thm:constrained}
If $a \neq b$, then there is no constrained world realizing the
assignment $\{a \mapsto \top, b \mapsto \bot\}$.
\end{theorem}

\begin{lstlisting}
theorem constrained_independence_fails (a b : S) (hab : a ≠ b) :
    ¬ ∀ assignment : S → Prop,
      ∃ cw : ConstrainedWorld S a b,
        ∀ s, cw.val s ↔ assignment s := by
  intro h
  let bad : S → Prop := fun s => s = a
  obtain ⟨⟨w, hw⟩, hmatch⟩ := h bad
  have ha : w a := (hmatch a).mpr rfl
  have hb : w b := hw ha
  exact hab ((hmatch b).mp hb)
\end{lstlisting}

\paragraph{Weather model.}
A more concrete example: three states of affairs (rain, clouds,
snow) with the physical constraint \emph{rain implies clouds}.
The weather model restricts worlds to those satisfying this
constraint.

\begin{theorem}[Weather model]
\label{thm:weather}
In the weather model, there is no world where rain holds but
clouds does not.
\end{theorem}

\noindent
These results make explicit a point that is often asserted
informally: TLP~2.061 (``Atomic facts are independent of one
another'') is not a logical law but a constraint on the space of
possible worlds.  Our formalization turns this observation into a
proof.

% ================================================================
\section{Formal Limits: Expressibility and Equivalence}
\label{sec:limits}
% ================================================================

The deepest results of the formalization concern the boundaries
of what the object language can express.  These are not claims
within the Tractarian system but \emph{about} it---theorems at
the meta-level that formalize Wittgenstein's saying/showing
distinction.

% ----------------------------------------------------------------
\subsection{The \texttt{expresses} Relation}
\label{sec:expresses}
% ----------------------------------------------------------------

We introduce a typed bridge between the object language and the
metalanguage:

\begin{definition}[\texttt{expresses}]
\label{def:expresses}
A proposition $p : \mathsf{Proposition}\;\Sach$ \emph{expresses}
a meta-level property $P : \Prop$ if and only if $p$'s truth
value tracks $P$ in every world:
\[
  \expresses\;p\;P \;\;\stackrel{\mathrm{def}}{=}\;\;
  \forall\;w : \World\;\Sach,\;\; p.\eval\;w \leftrightarrow P
\]
\end{definition}

\begin{lstlisting}
def expresses (p : Proposition S) (P : Prop) : Prop :=
  ∀ w : World S, p.eval w ↔ P
\end{lstlisting}

\noindent
The type signature makes the object/meta distinction visible:
$p$ lives in the object language, $P$ lives in \Lean's
metalanguage, and \texttt{expresses} is the claim that they
correspond.  This is the formalization's key analytical device.

% ----------------------------------------------------------------
\subsection{Expressibility Collapse}
\label{sec:collapse}
% ----------------------------------------------------------------

\begin{theorem}[Expressibility collapse]
\label{thm:collapse}
If a proposition $p$ expresses a world-independent property
$P : \Prop$, then $p$ is a tautology or a contradiction.
\end{theorem}

\begin{proof}
Suppose $\expresses\;p\;P$ holds.  Then for any two worlds
$w_1, w_2$:
\[
  p.\eval\;w_1 \;\leftrightarrow\; P \;\leftrightarrow\;
  p.\eval\;w_2
\]
so $p$ has the same truth value in all worlds.  By excluded
middle on $p.\eval\;w_0$ for an arbitrary world~$w_0$, either
$p$ holds everywhere (tautology) or fails everywhere
(contradiction).
\end{proof}

\begin{lstlisting}
theorem saying_showing_triviality (q : Proposition S)
    (P : Prop) [Nonempty S]
    (h : expresses q P) :
    IsTautology q ∨ IsContradiction q := by
  apply world_constant_taut_or_contra
  intro w₁ w₂
  exact (h w₁).trans (h w₂).symm
\end{lstlisting}

\noindent
The contrapositive is immediate and philosophically sharper:

\begin{corollary}[Nontrivial propositions cannot express
world-independent content]
\label{cor:nontrivial}
If $p$ is nontrivial (i.e., true in some world and false in
another), then there is no $P : \Prop$ such that
$\expresses\;p\;P$.
\end{corollary}

\begin{lstlisting}
theorem nontrivial_cannot_express_world_independent
    (p : Proposition S) (P : Prop) [Nonempty S]
    (hnt : Nontrivial p) :
    ¬ expresses p P := by
  intro h
  exact nontrivial_expressibility_requires_world_dependence
    p P hnt h
\end{lstlisting}

\noindent
This is the paper's headline result.  It says: if a proposition
genuinely ``says'' something---i.e., carves out a nontrivial
region of logical space---then what it says cannot be a
world-independent truth.  Genuine content requires variation
across possible worlds.  What is world-independent can only be
``shown'' by the structure of the formalism itself, not ``said''
within it.

This transforms the saying/showing distinction from a
philosophical thesis into a provable constraint on
truth-functional languages.

% ----------------------------------------------------------------
\subsection{The Equivalence Hierarchy}
\label{sec:hierarchy}
% ----------------------------------------------------------------

The formalization distinguishes three notions of proposition
equivalence, ordered by strictness:

\begin{definition}[Equivalence hierarchy]
\label{def:hierarchy}
Let $p, q : \mathsf{Proposition}\;\Sach$.
\begin{enumerate}[nosep]
  \item \textbf{Structural equivalence} ($p \structEq q$):
    $p$ and $q$ are built from identical constructors in the
    same tree shape with identical atoms.
  \item \textbf{Logical-form equivalence} ($p \formEq q$):
    there exists a bijective renaming of atoms (a permutation
    $\sigma : \Sach \xrightarrow{\sim} \Sach$) such that
    $p.\mathtt{rename}\;\sigma = q$.
  \item \textbf{Semantic equivalence} ($p \semEq q$):
    $p$ and $q$ have the same truth value in every world:
    $\forall w,\; p.\eval\;w \leftrightarrow q.\eval\;w$.
\end{enumerate}
\end{definition}

\noindent
We prove that $\formEq$ is an equivalence relation (reflexive by
the identity permutation, symmetric by inverse, transitive by
composition) and establish the strict hierarchy:

\begin{theorem}[Strict separation]
\label{thm:hierarchy}
$\structEq \;\subsetneq\; \formEq \;\subsetneq\; \semEq$
\end{theorem}

\begin{proof}
Each inclusion is witnessed by containment and a separating example.

\emph{$\structEq \subsetneq \formEq$:}
Structural equality implies propositional equality (by induction),
which implies logical-form equivalence via the identity
permutation.  Strictness: $\mathtt{elementary}(\mathit{rain})$
and $\mathtt{elementary}(\mathit{snow})$ are $\formEq$ (via the
swap permutation) but not $\structEq$ (different atoms).

\emph{$\formEq \subsetneq \semEq$:}
Logical-form equivalence implies truth-table isomorphism
(Theorem~\texttt{formEq\_implies\_truth\_table\_iso}), and
truth-table isomorphism with trivial relabeling gives semantic
equivalence.  Strictness: $\lnot\lnot(\mathtt{elementary}\;s)$
is $\semEq$ to $\mathtt{elementary}\;s$ (by double negation) but
not $\formEq$ (renaming preserves tree depth, so a depth-3 tree
cannot become depth-1).
\end{proof}

\noindent
This hierarchy clarifies a long-standing ambiguity in the
\emph{Tractatus}.  Wittgenstein's notion of ``logical form''
(TLP~4.0141) is neither syntactic identity nor truth-conditional
equivalence but something in between.  Our $\formEq$---defined
via bijective atom renaming---provides a formal candidate:
it captures the connective structure while abstracting over
which particular atoms appear.  The strict separation proves
that this intermediate notion is genuine, not collapsible into
either extreme.

% ----------------------------------------------------------------
\subsection{First-Order Extension}
\label{sec:firstorder}
% ----------------------------------------------------------------

In a companion file, we extend the propositional calculus with
first-order quantifiers using higher-order abstract syntax
(HOAS):

\begin{lstlisting}
inductive FOProp (S : Type) (D : Type) : Type where
  | lift     : Proposition S → FOProp S D
  | negFO    : FOProp S D → FOProp S D
  | conjFO   : FOProp S D → FOProp S D → FOProp S D
  | forall_  : (D → FOProp S D) → FOProp S D
  | exists_  : (D → FOProp S D) → FOProp S D
\end{lstlisting}

\noindent
The HOAS encoding leverages \Lean's metalanguage for variable
binding, avoiding explicit substitution machinery.  We prove
that compositionality lifts to the first-order level:

\begin{theorem}[First-order compositionality]
\label{thm:fo-comp}
For any $p : \mathsf{FOProp}\;\Sach\;D$ and worlds $w_1, w_2$
agreeing on all states of affairs, $p.\eval\;w_1 \leftrightarrow
p.\eval\;w_2$.
\end{theorem}

\noindent
We also establish universal--existential duality (TLP~5.52),
quantifier distribution, vacuous quantifier elimination, and De
Morgan duality for quantifiers.  The file notes the
Geach--Soames critique~\cite{geach1981,soames1983}: for infinite
domains, Wittgenstein's claim that quantifiers reduce to
applications of the $N$-operator is problematic, since the
$N$-operator is a finitary syntactic operation.

% ================================================================
\section{Related Work}
\label{sec:related}
% ================================================================

\paragraph{Formal reconstructions of the \emph{Tractatus}.}
Lokhorst~\cite{lokhorst1988} provides the most comprehensive
prior formal reconstruction, covering ontology, semantics, and
propositional attitudes in a set-theoretic framework.  His
axiomatization of independence and his truth-functionality
theorem (every sentence is a truth-function of elementary
sentences) overlap with our encoding, but he does not formalize
the saying/showing boundary itself.
Weiss~\cite{weiss2017} gives a 50-page reconstruction of the
\emph{Tractatus}'s logical system, analyzing the $N$-operator
and form-series in detail and proving that tautology-hood is
$\Pi^1_1$-complete for countable domains---a striking
complexity result that sharpens the Geach--Soames critique.
Miller~\cite{miller1995} develops a Tractarian semantics for
predicate logic that is independent of the metaphysics of
logical atomism.
Stokhof~\cite{stokhof2008} reads the \emph{Tractatus} as a
proto-formal semantics, tracing its influence through Carnap's
state descriptions to possible-worlds semantics, and argues
that the saying/showing distinction arises from the
universalism of the framework---a philosophical diagnosis our
expressibility collapse theorem makes precise.
All of these works provide formal clarity but are carried out
on paper; none offers machine-checked proofs, and none
parameterizes over world models to isolate which claims are
structural invariants versus modeling assumptions.

\paragraph{The Geach--Soames critique.}
Geach~\cite{geach1981} argues that the $N$-operator cannot
ground quantification over infinite domains, since it is a
finitary operation.  Soames~\cite{soames1983} sharpens this:
the \emph{Tractatus}'s claim that all of logic reduces to
truth-functional operations on elementary propositions is
undermined when the domain of quantification is infinite.
Our formalization takes this critique seriously: the first-order
extension (Section~\ref{sec:firstorder}) uses standard
quantifiers rather than the $N$-operator, and notes the
finite-domain connection as a separate theorem to be proved only
when an $N$-operator formalization is available.

\paragraph{Proof assistants and philosophy.}
Benzm\"uller and colleagues~\cite{benzmuller2017} have used
Isabelle to formalize ontological arguments and modal logic
via shallow semantic embeddings of higher-order modal logic.
Fitelson and Zalta~\cite{fitelson2007} pioneer ``computational
metaphysics,'' using automated reasoners to verify theorems
in abstract-object theory.
Scherf~\cite{scherf2025} provides the most directly comparable
work: a machine-checked axiomatization of Advaita Ved\=anta in
\Lean~4 (69 axioms, 40+ theorems), demonstrating that complete
philosophical systems can be formalized in dependent type
theory.
Corfield~\cite{corfield2020} argues at book length that
homotopy type theory is a better logic for philosophy than
predicate logic, but provides no machine-checked formalizations.
To our knowledge, no prior work formalizes any part of the
\emph{Tractatus} in any proof assistant, nor uses the type
system to isolate the saying/showing boundary as a provable
constraint.

% ================================================================
\section{Discussion}
\label{sec:discussion}
% ================================================================

% ----------------------------------------------------------------
\subsection{What the Formalization Shows}
\label{sec:shows}
% ----------------------------------------------------------------

The formalization yields three kinds of insight.

First, it makes the \emph{Tractatus}'s implicit modeling
decisions explicit.  The choice $\World\;\Sach = \Sach \to
\Prop$ is not neutral: it builds independence into the
foundations.  By parameterizing over \texttt{WorldModel}, we
separate what follows from the propositional syntax (compositionality)
from what follows from the choice of world space (independence).
This distinction is available in principle but rarely made
precise in the philosophical literature.

Second, it transforms the saying/showing distinction from a
metaphilosophical thesis into a theorem with a two-line proof.
The \texttt{expresses} relation is the key: once the
object/meta boundary is a type distinction, the collapse result
(Theorem~\ref{thm:collapse}) follows by transitivity of
biconditional.  The philosophical content lies not in the proof's
difficulty but in the \emph{formulation}: the act of making
\texttt{expresses} a typed relation reveals a structural
constraint that is invisible when the object/meta boundary is
handled informally.

Third, the equivalence hierarchy (Theorem~\ref{thm:hierarchy})
provides a precise formal candidate for Wittgenstein's
notion of ``logical form.''  The strict separation $\structEq
\subsetneq \formEq \subsetneq \semEq$ shows that logical form,
understood as connective structure up to atom permutation, is a
genuine intermediate notion that cannot be reduced to either
syntax or semantics.

% ----------------------------------------------------------------
\subsection{Limitations}
\label{sec:limitations}
% ----------------------------------------------------------------

Several aspects of the \emph{Tractatus} lie beyond the scope of
this formalization.

\paragraph{Object structure.}
The type $\Sach$ is abstract: we do not model how objects combine
into states of affairs.  TLP~2.0141 (``The possibility of its
occurrence in atomic facts is the form of the object'') suggests
a richer dependent-type encoding where the form of an object
constrains its combinatorial possibilities.  This would make
independence a substantive theorem rather than a definitional
consequence.

\paragraph{The $N$-operator.}
Wittgenstein's claim that all propositions can be generated from
elementary propositions by successive applications of a single
operation (TLP~6) is not formalized here.  The $N$-operator is
a finitary syntactic construction, and its connection to
quantification over infinite domains remains philosophically
contentious (Section~\ref{sec:firstorder}).

\paragraph{The picture theory.}
The \emph{Tractatus}'s claim that propositions are
\emph{pictures} of reality (TLP~2.1--2.174) involves a notion
of structural isomorphism between language and world that goes
beyond truth-functional evaluation.  Our formalization captures
the truth-conditional aspect but not the pictorial-isomorphism
aspect.

\paragraph{Self-reference and the ladder.}
TLP~6.54 famously declares that the propositions of the
\emph{Tractatus} itself are ``nonsensical'' and must be ``thrown
away'' after use.  Our formalization is, by design, a
metalanguage that talks \emph{about} the Tractarian object
language.  It cannot formalize its own status as a ladder to be
discarded---this is the point where type theory, like any formal
system, reaches its own saying/showing boundary.

We mark this boundary with a deliberate axiom:

\begin{lstlisting}
axiom silence : True
\end{lstlisting}

\noindent
This is not a missing proof but a design boundary.  The axiom
states nothing contentful ($\mathtt{True}$ is already provable);
it marks the point where the formalization acknowledges its own
limits.

% ================================================================
\section{Conclusion}
\label{sec:conclusion}
% ================================================================

We have presented a machine-checked reconstruction of the
semantic core of \emph{Tractatus Logico-Philosophicus} in
\Lean~4.  The formalization proves 46 theorems, decomposes the
\emph{Tractatus} into structural invariants, model-dependent
assumptions, and formal limits of expressibility, and
characterizes the saying/showing distinction as a theorem:
nontrivial propositions cannot express world-independent truths.

The equivalence hierarchy $\structEq \subsetneq \formEq
\subsetneq \semEq$ provides a typed decomposition of the space
between syntax and semantics, isolating logical form as a
genuine intermediate structure.

These results suggest that proof assistants are a productive tool
for philosophical analysis---not merely for verifying known
results, but for revealing structural constraints that are
invisible when the analysis is conducted on paper.

The formalization is publicly available at
\url{https://leangenius.org/proof/tractatus-ontology}.

% ================================================================
% References
% ================================================================

\begin{thebibliography}{99}

\bibitem{wittgenstein1921}
L.~Wittgenstein,
\emph{Tractatus Logico-Philosophicus},
Kegan Paul, 1921.
Translated by C.~K. Ogden; reprinted by Routledge, 2001.

\bibitem{lokhorst1988}
G.-J.~C. Lokhorst,
``Ontology, semantics and philosophy of mind in Wittgenstein's
\emph{Tractatus}: A formal reconstruction,''
\emph{Erkenntnis}, vol.~29, no.~1, pp.~35--75, 1988.
doi:~\texttt{10.1007/BF00166365}.

\bibitem{stokhof2008}
M.~Stokhof,
``The architecture of meaning: Wittgenstein's \emph{Tractatus} and
formal semantics,''
in \emph{Wittgenstein's Enduring Arguments},
D.~Levy and E.~Zamuner, Eds.\
London: Routledge, 2008, pp.~211--244.

\bibitem{weiss2017}
M.~Weiss,
``Logic in the \emph{Tractatus},''
\emph{The Review of Symbolic Logic}, vol.~10, no.~1,
pp.~1--50, 2017.
doi:~\texttt{10.1017/S1755020316000472}.

\bibitem{geach1981}
P.~Geach,
``Wittgenstein's operator $N$,''
\emph{Analysis}, vol.~41, no.~4, pp.~168--171, 1981.

\bibitem{soames1983}
S.~Soames,
``Generality, truth functions, and expressive capacity in the
\emph{Tractatus},''
\emph{The Philosophical Review}, vol.~92, no.~4, pp.~573--589, 1983.

\bibitem{miller1995}
H.~Miller,
``Tractarian semantics for predicate logic,''
\emph{History and Philosophy of Logic}, vol.~16, no.~2,
pp.~197--215, 1995.

\bibitem{benzmuller2017}
C.~Benzm\"uller and B.~Woltzenlogel Paleo,
``An object-logic explanation of the inconsistency in G\"odel's
ontological theory,''
\emph{Journal of Applied Logic}, vol.~24, pp.~253--265, 2017.

\bibitem{fitelson2007}
B.~Fitelson and E.~N. Zalta,
``Steps toward a computational metaphysics,''
\emph{Journal of Philosophical Logic}, vol.~36, no.~2,
pp.~227--247, 2007.

\bibitem{scherf2025}
M.~Scherf,
``The formal logic of Advaita Ved\=anta: A complete
axiomatization and consistency proof,''
PhilPapers, 2025.
Available: \url{https://github.com/matthew-scherf/Advaita}.

\bibitem{corfield2020}
D.~Corfield,
\emph{Modal Homotopy Type Theory: The Prospect of a New Logic
for Philosophy},
Oxford University Press, 2020.

\bibitem{moura2021}
L.~de Moura and S.~Ullrich,
``The \textsc{Lean}~4 theorem prover and programming language,''
in \emph{Proceedings of CADE-28}, Springer LNCS, 2021.

\end{thebibliography}

\end{document}
